% PAQUETES

% Paquetes de idioma y codificación principales
\usepackage[utf8]{inputenc} % Permite usar caracteres UTF-8.
\usepackage[T1]{fontenc} % Mejora la codificación de las fuentes.
\usepackage[spanish]{babel} % Adapta el documento al idioma español.
\usepackage{csquotes} % Gestiona correctamente las comillas
\usepackage{subcaption} % Permite subfiguras y subtítulos en figuras.
\usepackage{graphicx} % Permite insertar y ajustar imágenes.
\usepackage{placeins} % Controla la posición de las figuras y tablas.
\usepackage{float} % Permite fijar las figuras en la posición indicada.
\usepackage{booktabs} % Crea tablas con formato profesional.
\usepackage{bookmark} % Mejora la generación de marcadores PDF.
\usepackage{amsmath} % Añade herramientas para fórmulas matemáticas.
\usepackage[
    backend=biber,
    style=apa
]{biblatex} % Gestiona las referencias bibliográficas en formato APA.

% Paquetes de diseño y formato
\usepackage{xcolor} % Permite definir y usar colores personalizados.
\usepackage{tikz} % Permite crear gráficos y formas vectoriales.
\usepackage{tocloft}[v=2]

% Carga el archivo de referencias bibliográficas.
\addbibresource{bibliography/references.bib}